\chapter{Bẫy Ảo Danh Và Phô Diễn}
\markboth{Bẫy Ảo Danh Và Phô Diễn}{The Trap of Vanity and Showing Off}

\section{Muốn Được Biết Tên Nhanh Hơn Muốn Làm Giỏi}
\begin{truyen}{Muốn Được Biết Tên Nhanh Hơn Muốn Làm Giỏi}{Wanting Recognition Faster Than Mastery}
\chuhoa{C}{ó người trẻ rất khát cảm giác được nhắc tên, được biết tới, được chú ý, dù phần năng lực bên dưới còn mỏng.}
\textit{(Some young people crave being named, known, and noticed even while the ability underneath remains thin.)}

Sự nổi lên nhanh cho cảm giác mình đang đi đúng hướng.
\textit{(Quick attention creates the feeling of heading in the right direction.)}

Nhưng nếu danh đi trước thực lực quá xa, người ta sẽ phải dùng thêm rất nhiều sức để giữ một hình ảnh chưa đủ nền.
\textit{(But if recognition runs too far ahead of real ability, a person must spend enormous energy holding up an image with too little foundation.)}

Được biết đến là một việc phụ. Làm chắc mới là việc chính.
\textit{(Being known is secondary. Becoming solid is primary.)}

\end{truyen}

\begin{baihoc}
Danh tiếng đi trước năng lực quá xa thường biến thành áp lực phải diễn.
\textit{(Fame that runs too far ahead of ability often turns into pressure to perform.)}
\end{baihoc}

\begin{ghinhoanh}
\textbf{Short English to remember:}

Fame that runs too far ahead of ability often turns into pressure to perform.

\end{ghinhoanh}

\begin{tuhocnhanh}
\textbf{Câu hỏi tự học / Self-study questions}

1. Nhân vật đang muốn điều gì trước? \textit{What does the character want first?}

2. Vì sao khoảng cách này nguy hiểm? \textit{Why is this gap dangerous?}

3. Em nên xây nền nào trước? \textit{What foundation should you build first?}
\end{tuhocnhanh}
\ngancach

\section{Đăng Thành Tựu Nhiều Hơn Làm Thành Tựu}
\begin{truyen}{Đăng Thành Tựu Nhiều Hơn Làm Thành Tựu}{Posting Achievement More Than Building It}
\chuhoa{C}{ó người rất chăm kể tiến độ, khoe kế hoạch, công bố mục tiêu, nhưng phần làm thật lại mỏng và đứt quãng.}
\textit{(Some people are very active in sharing progress, displaying plans, and announcing goals while the real work remains thin and inconsistent.)}

Mỗi lần được phản ứng tốt, họ thấy như mình đã tiến thêm một đoạn.
\textit{(Each good reaction makes them feel as if they have already moved forward.)}

Nhưng cảm giác được ghi nhận sớm đôi khi làm động lực cho phần việc âm thầm yếu đi.
\textit{(But early validation can weaken the motivation for the quiet work itself.)}

Nếu nói về hành trình nhiều hơn đi đường, mình rất dễ yêu hình ảnh người đang cố gắng hơn là thật sự cố gắng.
\textit{(If you talk about the journey more than you walk it, you may fall in love with the image of effort rather than effort itself.)}

\end{truyen}

\begin{baihoc}
Đừng để sự công bố thay thế cho tiến độ thật.
\textit{(Do not let public announcement replace real progress.)}
\end{baihoc}

\begin{ghinhoanh}
\textbf{Short English to remember:}

Do not let public announcement replace real progress.

\end{ghinhoanh}

\begin{tuhocnhanh}
\textbf{Câu hỏi tự học / Self-study questions}

1. Nhân vật đang nuôi điều gì bằng phản ứng bên ngoài? \textit{What is external reaction feeding?}

2. Điều gì đang yếu đi? \textit{What is getting weaker?}

3. Em có cần bớt công bố để làm sâu hơn không? \textit{Do you need less announcement and more depth?}
\end{tuhocnhanh}
\ngancach

\section{Phô Diễn Điều Mình Chưa Giữ Được}
\begin{truyen}{Phô Diễn Điều Mình Chưa Giữ Được}{Showing Off What You Cannot Sustain}
\chuhoa{C}{ó người thích phô bày mức sống, quan hệ, thành công, hoặc độ trưởng thành mà mình thật ra chưa giữ được bền.}
\textit{(Some people like to display a lifestyle, network, success, or maturity they cannot actually sustain.)}

Sự phô diễn cho họ vài phút được nhìn lên.
\textit{(Showing off gives them a few moments of being looked up to.)}

Nhưng sau đó là áp lực phải tiếp tục giữ vai, tiếp tục tiêu, tiếp tục tỏ ra đủ.
\textit{(But afterward comes the pressure to keep playing the role, keep spending, and keep appearing sufficient.)}

Càng khoe thứ chưa bền, mình càng sống trong sợ hãi bị lộ khoảng trống.
\textit{(The more you display what is not stable, the more you live in fear of being exposed.)}

\end{truyen}

\begin{baihoc}
Thứ gì chưa đủ nền thì chưa cần mang ra làm sân khấu.
\textit{(What lacks foundation does not need a stage yet.)}
\end{baihoc}

\begin{ghinhoanh}
\textbf{Short English to remember:}

What lacks foundation does not need a stage yet.

\end{ghinhoanh}

\begin{tuhocnhanh}
\textbf{Câu hỏi tự học / Self-study questions}

1. Nhân vật đang phô điều gì? \textit{What is the character displaying?}

2. Điều gì đến sau vài phút được chú ý? \textit{What comes after a few moments of attention?}

3. Em có đang đưa ra ngoài thứ còn quá non không? \textit{Are you putting something outside before it is ready?}
\end{tuhocnhanh}
\ngancach

\section{Lấy Ánh Nhìn Bên Ngoài Làm Nhiên Liệu Sống}
\begin{truyen}{Lấy Ánh Nhìn Bên Ngoài Làm Nhiên Liệu Sống}{Using External Attention as Emotional Fuel}
\chuhoa{N}{hiều người chỉ thấy mình hăng khi được chú ý, được nhắc tới, được khen ngợi, hoặc được ngưỡng mộ.}
\textit{(Many people feel energized only when they are noticed, mentioned, praised, or admired.)}

Khi ánh nhìn đó giảm đi, họ xuống tinh thần rất nhanh.
\textit{(When that attention fades, their spirit drops quickly.)}

Một đời sống như vậy rất khó yên vì động cơ nằm hết trong tay người khác.
\textit{(A life like that struggles to stay steady because its fuel source lies in other people’s hands.)}

Muốn đi xa, một người trẻ phải học cách làm việc ngay cả khi không ai nhìn.
\textit{(To go far, a young person must learn to work even when no one is watching.)}

\end{truyen}

\begin{baihoc}
Sự bền bỉ thật được nuôi bằng ý thức bên trong, không chỉ bằng ánh đèn bên ngoài.
\textit{(Real persistence is fed by inner conviction, not only by external spotlight.)}
\end{baihoc}

\begin{ghinhoanh}
\textbf{Short English to remember:}

Real persistence is fed by inner conviction, not only by external spotlight.

\end{ghinhoanh}

\begin{tuhocnhanh}
\textbf{Câu hỏi tự học / Self-study questions}

1. Nhiên liệu của nhân vật đang đến từ đâu? \textit{Where is the character’s fuel coming from?}

2. Vì sao kiểu nhiên liệu này bấp bênh? \textit{Why is this fuel unstable?}

3. Em làm được điều gì ngay cả khi không ai nhìn? \textit{What can you keep doing even when no one is watching?}
\end{tuhocnhanh}
\ngancach

\section{Nói Về Hành Trình Như Người Đã Tới Đích}
\begin{truyen}{Nói Về Hành Trình Như Người Đã Tới Đích}{Talking About the Journey as if You Already Arrived}
\chuhoa{C}{ó người mới đi được một đoạn ngắn đã bắt đầu kể chuyện như thể mình đã thành hình rất vững.}
\textit{(Some people have only gone a short distance yet already speak as if they have fully arrived.)}

Cách kể đó đem lại sự chú ý nhanh và cảm giác mình có vị thế.
\textit{(That kind of storytelling brings attention quickly and creates a sense of status.)}

Nhưng nếu lời kể chạy trước chặng đường thật quá xa, con người sẽ dần sống để giữ hình ảnh hơn là tiếp tục lớn lên.
\textit{(But if the story runs too far ahead of the real journey, a person begins living to maintain an image rather than keep growing.)}

Ảo danh thường bắt đầu ở chỗ câu chuyện về mình đẹp hơn sự thật một khoảng vừa đủ để người khác ngưỡng mộ.
\textit{(False glory often begins where the story about you becomes just beautiful enough to exceed the truth.)}

\end{truyen}

\begin{baihoc}
Hành trình thật không cần kể quá tay để có giá trị.
\textit{(A real journey does not need exaggeration to have value.)}
\end{baihoc}

\begin{ghinhoanh}
\textbf{Short English to remember:}

A real journey does not need exaggeration to matter.

\end{ghinhoanh}

\begin{tuhocnhanh}
\textbf{Câu hỏi tự học / Self-study questions}

1. Nhân vật đang kể quá phần nào? \textit{What part is the character overstating?}

2. Việc này đẩy họ sống vì điều gì? \textit{What does this push them to live for?}

3. Em có đang kể mình đẹp hơn sự thật không? \textit{Are you presenting yourself more beautifully than the truth?}
\end{tuhocnhanh}
\ngancach

\section{Muốn Nổi Trong Mắt Người Lạ Hơn Có Ích Cho Người Gần}
\begin{truyen}{Muốn Nổi Trong Mắt Người Lạ Hơn Có Ích Cho Người Gần}{Wanting to Impress Strangers More Than Helping Those Close to You}
\chuhoa{C}{ó người rất đầu tư cho cách mình xuất hiện trước đám đông, nhưng lại sơ sài với những bổn phận gần nhất quanh mình.}
\textit{(Some people invest heavily in how they appear before the public, yet stay careless with the responsibilities closest to them.)}

Họ thích tiếng vang ở xa hơn là giá trị thật ở gần.
\textit{(They prefer distant applause over nearby real value.)}

Nhưng đời sống đáng tin thường được đo nhiều hơn ở cách em hiện diện với người gần, không chỉ ở cách em được nhìn từ xa.
\textit{(But a trustworthy life is measured more by how you show up for people nearby than by how you look from afar.)}

Phô diễn dễ làm người lạ chú ý. Hữu ích mới khiến đời sống có trọng lượng.
\textit{(Display easily attracts strangers. Usefulness gives a life real weight.)}

\end{truyen}

\begin{baihoc}
Đừng để ánh nhìn xa làm em quên trách nhiệm gần.
\textit{(Do not let distant attention make you forget nearby responsibilities.)}
\end{baihoc}

\begin{ghinhoanh}
\textbf{Short English to remember:}

Distant attention must not replace nearby responsibility.

\end{ghinhoanh}

\begin{tuhocnhanh}
\textbf{Câu hỏi tự học / Self-study questions}

1. Nhân vật đang đầu tư cho điều gì? \textit{What is the character investing in?}

2. Điều gì đang bị sơ sài? \textit{What is being neglected?}

3. Em đang muốn nổi hơn hay muốn hữu ích hơn? \textit{Do you want to stand out more, or to be more useful?}
\end{tuhocnhanh}
\ngancach

\section{Sửa Đời Theo Góc Máy Hơn Theo Giá Trị}
\begin{truyen}{Sửa Đời Theo Góc Máy Hơn Theo Giá Trị}{Shaping Life for the Camera More Than for Values}
\chuhoa{C}{ó người dần lựa chọn nơi đi, điều làm, và cả cách sống theo việc nó có đẹp để đưa ra ngoài hay không.}
\textit{(Some people slowly choose places, actions, and even lifestyles by whether they look good for display.)}

Cuộc sống của họ trở nên rất biết tạo khung hình.
\textit{(Their life becomes very skilled at creating a frame.)}

Nhưng một đời sống hợp góc máy chưa chắc là một đời sống hợp giá trị.
\textit{(But a life that fits the camera is not necessarily a life that fits real values.)}

Khi khung hình được ưu tiên hơn nội dung, con người rất dễ sống ngày càng xa điều thật sự nuôi mình.
\textit{(When framing outranks substance, people easily drift away from what truly nourishes them.)}

\end{truyen}

\begin{baihoc}
Đẹp để đăng không quý bằng đúng để sống.
\textit{(Looking good to post is less valuable than living rightly.)}
\end{baihoc}

\begin{ghinhoanh}
\textbf{Short English to remember:}

Looking good to post matters less than living rightly.

\end{ghinhoanh}

\begin{tuhocnhanh}
\textbf{Câu hỏi tự học / Self-study questions}

1. Nhân vật đang ưu tiên điều gì? \textit{What is the character prioritizing?}

2. Điều gì bị đẩy xuống sau? \textit{What is being pushed aside?}

3. Em có đang chọn theo góc nhìn hơn là theo giá trị không? \textit{Are you choosing based on appearance more than value?}
\end{tuhocnhanh}
\ngancach

\section{Coi Im Lặng Là Thất Bại}
\begin{truyen}{Coi Im Lặng Là Thất Bại}{Treating Silence as Failure}
\chuhoa{C}{ó người hễ không còn ai nhắc tới mình, không còn tương tác, không còn ánh nhìn là thấy như mình đang tụt xuống.}
\textit{(Some people feel that whenever nobody mentions them, interacts, or looks at them, they must be declining.)}

Sự im lặng làm họ hoảng và vội tạo thêm tiếng động quanh mình.
\textit{(Silence makes them panic, so they rush to create more noise around themselves.)}

Nhưng có những giai đoạn im lặng rất cần cho việc làm dày kỹ năng, làm chắc tính cách, và làm sâu công việc.
\textit{(But some silent phases are necessary for deepening skill, strengthening character, and maturing work.)}

Ai coi im lặng là thất bại sẽ rất khó đi đủ xa để có chiều sâu thật.
\textit{(Anyone who treats silence as failure will struggle to go far enough to build real depth.)}

\end{truyen}

\begin{baihoc}
Không phải lúc ít được chú ý là lúc em ít giá trị.
\textit{(Being less noticed does not mean being less valuable.)}
\end{baihoc}

\begin{ghinhoanh}
\textbf{Short English to remember:}

Less attention does not mean less value.

\end{ghinhoanh}

\begin{tuhocnhanh}
\textbf{Câu hỏi tự học / Self-study questions}

1. Vì sao nhân vật sợ im lặng? \textit{Why is the character afraid of silence?}

2. Im lặng có thể giúp điều gì? \textit{What can silence help build?}

3. Em có chịu được giai đoạn không được chú ý không? \textit{Can you endure seasons of low attention?}
\end{tuhocnhanh}
\ngancach

\section{Phát Sáng Bên Ngoài, Rỗng Dần Bên Trong}
\begin{truyen}{Phát Sáng Bên Ngoài, Rỗng Dần Bên Trong}{Shining Outward While Growing Empty Inside}
\chuhoa{C}{ó người càng được nhìn thấy nhiều lại càng kiệt dần ở phần riêng tư không ai biết.}
\textit{(Some people become more visible outwardly while growing more exhausted in the private part no one sees.)}

Họ liên tục duy trì năng lượng cho bề mặt nên không còn chỗ cho việc nuôi phần sâu bên trong.
\textit{(They keep pouring energy into the surface and leave little room to nourish the inner life.)}

Một lúc nào đó, ánh sáng bên ngoài vẫn còn nhưng trong lòng đã gần như cạn.
\textit{(At some point, the outer shine remains while the inside is nearly empty.)}

Ảo danh nguy hiểm nhất là khi nó thắp sáng hình ảnh của em bằng chính phần sống thật đang bị rút cạn.
\textit{(The most dangerous false glory is the kind that lights your image using the fuel of your depleted real life.)}

\end{truyen}

\begin{baihoc}
Hình ảnh không đáng để đổi bằng phần sống thật bị rút khô.
\textit{(No image is worth drying out your real life for.)}
\end{baihoc}

\begin{ghinhoanh}
\textbf{Short English to remember:}

No image is worth drying out your real life for.

\end{ghinhoanh}

\begin{tuhocnhanh}
\textbf{Câu hỏi tự học / Self-study questions}

1. Điều gì đang phát sáng? \textit{What is shining?}

2. Điều gì đang rỗng dần? \textit{What is becoming empty?}

3. Em cần nuôi lại phần thật nào của mình? \textit{What real part of yourself needs to be nourished again?}
\end{tuhocnhanh}
\ngancach

\section{Làm Thật Lặng Lẽ Vẫn Là Một Kiểu Danh Dự}
\begin{truyen}{Làm Thật Lặng Lẽ Vẫn Là Một Kiểu Danh Dự}{Quiet Honest Work Is Also a Form of Honor}
\chuhoa{C}{ó người sau nhiều vòng mỏi mệt đã chọn bớt phô diễn, bớt đẩy hình ảnh, và quay lại làm việc chắc hơn trong yên lặng.}
\textit{(Some people, after many tiring cycles, choose less display, less image pushing, and return to steadier work in quiet.)}

Ban đầu họ thấy như mình mờ đi trong mắt đám đông.
\textit{(At first, they feel as if they are fading from the crowd's eyes.)}

Nhưng dần dần, họ có lại một sự vững chắc mà tiếng ồn trước đó không đem lại được.
\textit{(But gradually, they regain a solidity that the earlier noise never gave them.)}

Danh dự không chỉ nằm ở việc được thấy. Nó còn nằm ở việc làm điều đáng trọng ngay cả khi không có ai vỗ tay.
\textit{(Honor is not only in being seen. It is also in doing what is worthy even when no one applauds.)}

\end{truyen}

\begin{baihoc}
Im lặng làm thật nhiều khi đáng giá hơn ồn ào được thấy.
\textit{(Quiet honest work is often worth more than noisy visibility.)}
\end{baihoc}

\begin{ghinhoanh}
\textbf{Short English to remember:}

Quiet honest work is often worth more than noisy visibility.

\end{ghinhoanh}

\begin{tuhocnhanh}
\textbf{Câu hỏi tự học / Self-study questions}

1. Nhân vật đã chọn bỏ bớt điều gì? \textit{What has the character chosen to reduce?}

2. Điều gì được lấy lại? \textit{What is regained?}

3. Em có thể làm tốt điều gì ngay cả khi không ai vỗ tay? \textit{What can you do well even if nobody applauds?}
\end{tuhocnhanh}
